\documentclass[aspectratio=169]{beamer}
\usepackage[UTF8]{ctex}
\usepackage{booktabs}
\usepackage{array}
\usepackage{hyperref}

\usetheme{Madrid}
\usecolortheme{default}
\definecolor{brandgreen}{RGB}{29,191,115}
\definecolor{brandlight}{RGB}{244,246,243}
\setbeamercolor{palette primary}{bg=brandgreen,fg=white}
\setbeamercolor{palette secondary}{bg=brandgreen!80!black,fg=white}
\setbeamercolor{structure}{fg=brandgreen}
\setbeamercolor{block title}{bg=brandgreen!12,fg=black}
\setbeamercolor{block body}{bg=brandlight,fg=black}
\setbeamertemplate{navigation symbols}{}

\title{智慧课堂 Class Assistant}
\subtitle{软件工程探索与实践（2）课程项目汇报}
\author{颜宇翔 \and 魏天翼 \and 蒋天源}
\institute{在线课堂辅助学习桌面应用}
\date{2026年6月}

\begin{document}

\begin{frame}
  \titlepage
\end{frame}

\begin{frame}{项目背景}
  \begin{block}{使用场景}
    在线课堂里，学生往往需要同时处理课件阅读、课堂记录、重点理解、课后复习和课堂题目等任务。现有做法通常依赖浏览器、笔记工具和聊天工具来回切换，信息容易分散，上下文也很难连续保留。
  \end{block}

  \begin{block}{项目目标}
    我们希望做一个围绕真实课堂流程的桌面助手，把课件捕获、内容解析、笔记整理、AI 对话和课堂题目辅助整合在同一个应用里，让它既能服务上课过程，也能服务课后复盘。
  \end{block}
\end{frame}

\begin{frame}{整体方案}
  \begin{columns}[T]
    \begin{column}{0.48\textwidth}
      \begin{block}{桌面端}
        \begin{itemize}
          \item 使用 Electron 搭建统一桌面界面
          \item 通过 BrowserView 承载雨课堂页面
          \item 在同一窗口中组织课件、解析、笔记与对话
        \end{itemize}
      \end{block}

      \begin{block}{服务端}
        \begin{itemize}
          \item 基于 Node.js 与 Express 提供本地服务
          \item 通过 SSE 支持解析、对话、笔记的流式输出
          \item 本地维护捕获记录、课堂状态和笔记数据
        \end{itemize}
      \end{block}
    \end{column}
    \begin{column}{0.48\textwidth}
      \begin{block}{核心流程}
        \begin{itemize}
          \item 捕获课堂中的课件内容
          \item 对课件进行去重、OCR 和结构化解析
          \item 结合上下文完成问答、总结与笔记生成
          \item 在真实题目页面上执行课堂答题辅助
        \end{itemize}
      \end{block}

      \begin{block}{设计思路}
        \begin{itemize}
          \item 尽量贴近真实课堂使用环境
          \item 各模块共享课堂上下文，而不是各自独立工作
          \item 保留继续扩展更多课堂能力的接口
        \end{itemize}
      \end{block}
    \end{column}
  \end{columns}
\end{frame}

\begin{frame}{课件捕获与解析}
  \begin{block}{课件接入方式}
    应用将雨课堂页面直接嵌入桌面端，用户不需要频繁在浏览器和本地工具之间切换。这样一方面便于保留完整的课堂上下文，另一方面也让后续捕获、解析和交互都能在同一流程中完成。
  \end{block}

  \begin{block}{处理链路}
    \begin{itemize}
      \item 监听页面中的课件图像与可见内容
      \item 对捕获结果进行尺寸过滤与去重
      \item 结合 OCR 提取文字信息，增强后续理解能力
      \item 调用模型生成结构化解析结果，包括课件分类、知识点与题目答案说明
    \end{itemize}
  \end{block}
\end{frame}

\begin{frame}{统一上下文与 RAG}
  \begin{block}{为什么要做统一上下文}
    课堂学习不是对单页课件的孤立理解，而是一个连续过程。只看当前页面，模型很容易遗漏前后关系；如果能结合邻近课件、历史解析结果和已有笔记，输出就会更接近真实学习需要。
  \end{block}

  \begin{block}{实现方式}
    我们将当前课件、相邻课件、历史分析结果和笔记内容组织成统一上下文，再作为检索增强信息提供给多个能力模块。这样，分析、深度思考、AI 对话和 AI 笔记生成都能共享同一套课堂记忆。
  \end{block}

  \begin{block}{带来的效果}
    \begin{itemize}
      \item 输出不再只依赖当前一页内容
      \item 更适合连续追问、知识串联和复习总结
      \item 降低不同模块之间信息割裂的问题
    \end{itemize}
  \end{block}
\end{frame}

\begin{frame}{笔记系统设计}
  \small
  \begin{block}{笔记组织方式}
    \begin{itemize}
      \item 每张课件对应一份独立笔记
      \item 手动笔记与 AI 笔记分区展示
      \item 支持导入解析结果、导出整理内容与后续复习使用
    \end{itemize}
  \end{block}

  \begin{block}{本次实现重点}
    \begin{itemize}
      \item 调整笔记预览逻辑，使预览与当前课件状态保持一致
      \item 优化 AI 笔记生成区域的交互方式，支持拖拽调整显示空间
      \item 将统一上下文接入笔记生成过程，使输出更贴合当前课堂内容
    \end{itemize}
  \end{block}
\end{frame}

\begin{frame}{GUI Agent 自动答题}
  \begin{block}{设计目标}
    这一部分服务的是课堂中的真实题目页面。系统先根据课件和题目内容生成结构化答案，再由 GUI Agent 在页面中识别元素并执行对应操作，从而完成选择、填空或主观题输入。
  \end{block}

  \begin{block}{工作流程}
    \begin{itemize}
      \item 课件捕获模块识别题目页面
      \item 解析结果中生成可提交的答案与说明
      \item GUI Agent 读取页面状态并选择操作对象
      \item 最终动作仍然执行在真实课堂页面中
    \end{itemize}
  \end{block}

  \begin{block}{系统意义}
    这部分说明应用不是离线分析器，而是一个能参与真实课堂交互流程的学习辅助工具。
  \end{block}
\end{frame}

\begin{frame}{其他模块}
  \begin{block}{OCR}
    OCR 用于补充图像课件中的文字信息。当页面中存在较多图片文字、题干文本或局部说明时，OCR 能为后续解析和问答提供更完整的输入基础。
  \end{block}

  \begin{block}{前端交互}
    前端围绕课堂场景做了统一组织，将课件区域、分析结果、笔记系统和对话面板放在同一界面中，降低切换成本，也方便用户围绕当前课件连续操作。
  \end{block}

  \begin{block}{程序打包}
    项目最终以桌面应用形式交付，便于直接在本地课堂环境中运行，而不是依赖额外的开发环境来启动全部功能。
  \end{block}
\end{frame}

\begin{frame}{团队分工}
  \small
  \begin{table}
    \centering
    \renewcommand{\arraystretch}{1.2}
    \begin{tabular}{>{\raggedright\arraybackslash}p{0.16\textwidth} >{\raggedright\arraybackslash}p{0.76\textwidth}}
      \toprule
      成员 & 主要分工 \\
      \midrule
      颜宇翔 & 主要框架与 RAG 相关工作，负责前端组织、程序打包、OCR 等除笔记系统和 GUI Agent 自动答题外的整体模块串联。 \\
      魏天翼 & 笔记系统与相关框架设计，负责笔记组织、预览、编辑与 AI 笔记能力的实现。 \\
      蒋天源 & GUI Agent 自动答题与相关框架设计，负责真实课堂题目页面下的识别、执行与交互链路。 \\
      \bottomrule
    \end{tabular}
  \end{table}
\end{frame}

\begin{frame}{总结}
  \begin{block}{项目完成情况}
    项目已经打通了在线课堂中的课件捕获、内容解析、上下文组织、笔记生成、对话辅助和课堂题目交互等核心链路，形成了一个较完整的课堂学习辅助桌面应用。
  \end{block}

  \begin{block}{后续工作}
    \begin{itemize}
      \item 继续增强课堂历史内容的组织与利用能力
      \item 提升笔记导出、复盘整理和长期学习支持
      \item 继续扩展 GUI Agent 在更多课堂页面场景下的适配能力
    \end{itemize}
  \end{block}

  \begin{center}
    \Large 谢谢老师
  \end{center}
\end{frame}

\end{document}
